\subsection{Time-Aligned Fully-Real Validation}
\label{sec:evaluation:realboth}

The strongest empirical validation in the paper combines the real Azure
Functions 2021 per-invocation trace with real ElectricityMaps
carbon-intensity data from the \emph{same calendar window} as the
trace (12--25 January 2021). All other experiments in
\S\ref{sec:evaluation} use synthetic-but-schema-faithful workload
calibrated to Azure characterizations \citep{shahrad2020serverless,
zhang2021faster,azuretrace}, or real carbon with synthetic workload,
or real workload with carbon from a later year. This section
substitutes both axes with real time-aligned data and reports the
findings.

\paragraph{Data sources.}
\begin{itemize}
\item \textbf{Workload}: the
\texttt{AzureFunctionsInvocationTraceForTwoWeeksJan2021} trace from
\citet{zhang2021faster,azuretrace} (1.98M invocations over 14 days,
loaded as a 24h slice giving 80{,}990 invocations across 119 unique
functions; round-robin assignment to the 5 regions). The trace
provides exact per-invocation arrival timestamps and execution
durations.
\item \textbf{Carbon}: ElectricityMaps hourly carbon-intensity feeds
for the five regions over 12--25 January 2021 (DE 396.6 g mean, FR
89.4 g, GB 293.3 g, US-CAISO 319.6 g, PL 801.1 g; range from FR's
diurnal flat 73--108 g to PL's coal-base-load 690--869 g, giving a
carbon-ratio span of roughly $11\times$ across the topology).
\end{itemize}

\paragraph{Headline result on fully-real time-aligned data.}

\begin{table}[h]
\centering\small
\caption{Real Azure 2021 trace + real ElectricityMaps Jan 2021 carbon
data, 24h, 80{,}990 invocations, 5 regions.}
\label{tab:realboth_headline}
\begin{tabular}{lrrrr}
\toprule
Scheduler         & Carbon (g) & vs FIFO & SLA & Cold \\
\midrule
FIFO              & 231.33     &   0.0\%   & 0.39\% & 1.68\% \\
Wait-Awhile-Batch & 239.29     & $-$3.44\% & 0.39\% & 2.25\% \\
Lechowicz         & 242.91     & $-$5.00\% & 0.39\% & 2.50\% \\
Spatial           & 73.11      & \textbf{+68.39\%} & 0.39\% & 1.68\% \\
GreenFaaS-v1      & 78.33      & +66.14\%  & 0.64\% & 2.99\% \\
\GreenFaaS        & \textbf{73.22} & \textbf{+68.35\%} & 0.39\% & 1.68\% \\
\bottomrule
\end{tabular}
\end{table}

Four observations.

\textbf{(1) GreenFaaS--Spatial gap collapses to 0.04 percentage
points} on fully-real time-aligned data --- tighter than the
synthetic 1.1pp gap and tighter than the real-LWA-carbon 0.5pp gap.
On real workload + real carbon, GreenFaaS and Spatial are
empirically indistinguishable. The \emph{topology-robustness} claim
from \S\ref{sec:evaluation:topology} stands: GreenFaaS matches
Spatial when Spatial works, and as
\S\ref{sec:evaluation:realtopology} shows, falls back to FIFO
byte-for-byte when Spatial does not.

\textbf{(2) Wait-Awhile and Lechowicz failure modes are
\emph{quantitatively softer} on real workload than on synthetic.}
Wait-Awhile loses only 3.44\% on real workload + real carbon, vs
$-306\%$ on synthetic. Lechowicz loses 5.00\%, vs $-238\%$ on
synthetic. The qualitative finding (batch carbon-aware techniques
fail on FaaS) holds; the catastrophic magnitudes seen in the
synthetic experiments were partly artifacts of the synthetic
workload's extreme warm-pool affinity (synthetic FIFO baseline
cold-start rate 0.07\% vs real workload's 1.68\% --- real FaaS has a
much higher baseline cold-start rate, so deferral can't make it
proportionally worse). This is itself a finding worth reporting
honestly.

\textbf{(3) The v1 ablation is 2.21 percentage points worse on real
data} (66.14\% vs 68.35\%) with cold-start rate doubled (2.99\% vs
1.68\%) and SLA violation up 60\% (0.64\% vs 0.39\%). The
\S\ref{sec:tradeoff:idle} idle-energy correction matters even more
on real data than synthetic, because the bursty real arrival pattern
exercises the warm-pool more aggressively than the synthetic Poisson
generator.

\textbf{(4) The do-no-harm property reproduces perfectly on fully-
real data.} The topology sweep on time-aligned data (data in
\code{results/real\_workload\_topology.csv}) shows
\GreenFaaS{} matching FIFO byte-for-byte in single-region DE and
single-region CAISO scenarios (297.39 g and 234.83 g respectively,
identical to FIFO), while v1 loses 5.04\% on single-region DE and
6.52\% on single-region CAISO. The
\S\ref{sec:tradeoff:idle} contribution is what produces this
property on real workload, just as it does on synthetic.

\paragraph{Spatial versus temporal decision breakdown.}
To make precise where the savings come from, we instrumented every
\GreenFaaS{} decision on the full 5-region fully-real experiment. Of
the 80{,}990 invocations, \GreenFaaS{} routes $17.5\%$ to a non-home
(cleaner) region, defers $0\%$ temporally, and leaves the
remaining $82.5\%$ at home executing immediately (these are
interactive functions pinned to home, or invocations whose home
region is already the cleanest reachable option). The carbon savings
on this experiment are therefore \emph{entirely} spatial: the
temporal-deferral axis contributes nothing here, because the
ElectricityMaps carbon trace has hourly resolution while the
deferrable deadline budget is on the order of minutes, so no
within-deadline deferral ever beats immediate spatial routing. This
is consistent with the forecast-invariance finding of
\S\ref{sec:evaluation:forecast}: on real grid data at hourly
resolution, GreenFaaS's value is its cold-start-aware spatial gate
and its do-no-harm fallback, not temporal shifting. The temporal
machinery becomes load-bearing only when sub-hourly carbon variation
is both present and forecastable --- a regime our current real data
does not exercise, and which we flag as a limitation.

\paragraph{Topology summary.} On time-aligned fully-real data,
across the five topologies of
\S\ref{sec:evaluation:realtopology}. Note that the full-5-region row
below ($+74.33\%$ / $+74.10\%$) differs from
Table~\ref{tab:realboth_headline}'s $+68.39\%$ / $+68.35\%$: the
table reports a single fixed 24h slice (the first window, 80{,}990
invocations), whereas the topology sweep below re-runs each topology
configuration over a different 24h window with its own invocation
count and carbon profile, so the absolute reductions differ while the
\GreenFaaS{}--Spatial gap stays within half a percentage point in
both. The table is the single-configuration benchmark; the sweep
reports the cross-topology pattern.

\begin{center}
\small
\begin{tabular}{lrrr}
\toprule
Topology               & Spatial   & \GreenFaaS{}     & gap \\
\midrule
Full 5-region          & +74.33\%  & +74.10\%   & $-0.23$ \\
EU-only (FR/DE/GB/PL)  & +83.42\%  & +83.17\%   & $-0.25$ \\
Coal-belt (DE/GB/PL)   & +53.10\%  & +52.61\%   & $-0.49$ \\
Single-region DE       &   0.0\%   & 0.0\% (=FIFO) & 0.00 \\
Single-region CAISO    &   0.0\%   & 0.0\% (=FIFO) & 0.00 \\
\bottomrule
\end{tabular}
\end{center}

Same shape as the synthetic-workload-plus-real-LWA-carbon results
(\S\ref{sec:evaluation:realtopology}): GreenFaaS matches Spatial to
within half a percentage point across multi-region topologies on
fully-real data, including the coal-belt topology (where GreenFaaS
neither beats nor is beaten by Spatial beyond run-to-run variance),
and the do-no-harm property holds. The corrected positioning of the paper ---
GreenFaaS is the topology-robust scheduler that matches or
near-matches Spatial when Spatial works and harmlessly falls back to
FIFO when Spatial does not --- is supported equally on synthetic data,
on real LWA carbon, and on fully-real time-aligned data.

\paragraph{Window-shift variance characterization on real workload.}
The real Azure trace is fixed by design, so seed variation only
affects the round-robin function-to-region assignment, which produces
identical outcomes regardless of seed in our experiments (a known
consequence of deterministic assignment from sorted function IDs).
We characterize run-to-run variance on real workload by varying the
24-hour \emph{window} within the 14-day trace instead: we re-ran the
fully-real headline configuration on five non-overlapping 24h windows
starting at days $\{0, 1, 2, 3, 4\}$ of the trace, with matching 24h
slices of the EM Jan 2021 carbon trace.

\begin{table}[h]
\centering\small
\caption{Window-shift variance, fully-real time-aligned data. Mean
$\pm$ std across 5 non-overlapping 24h windows.}
\label{tab:window_variance}
\begin{tabular}{lr}
\toprule
Scheduler          & vs FIFO (5-window mean $\pm$ std) \\
\midrule
Wait-Awhile        & $-3.85\% \pm 1.99\%$ \\
Lechowicz          & $-5.56\% \pm 1.81\%$ \\
Spatial            & $+54.44\% \pm 13.57\%$ \\
GreenFaaS-v1       & $+51.49\% \pm 13.62\%$ \\
\GreenFaaS         & $+54.38\% \pm 13.57\%$ \\
\bottomrule
\end{tabular}
\end{table}

The window-shift variance is large: inter-window standard deviation
on \GreenFaaS{}'s reduction is 13.6 percentage points. This reflects
the natural variation in 24h workload intensity (windows ranging from
81k to 290k invocations) and per-window carbon profile (winter day-to-day
shifts in Polish coal vs French nuclear output). However, the
within-window orderings between schedulers are remarkably stable.

In other words: \emph{the absolute carbon footprint varies widely
window-to-window, but the scheduler ranking is stable to within
fractions of a percentage point.} Concretely:

\begin{itemize}
\item \GreenFaaS{}--Spatial gap mean: 0.058pp,
window-to-window std: 0.032pp. \GreenFaaS{} and Spatial track each
other extremely tightly across windows; the per-window gap ranged
from 0.022pp to 0.110pp, never out of $4\sigma$ of the mean.
\item \GreenFaaS{}--v1 gap mean: 2.89pp,
window-to-window std: 1.31pp. The \S\ref{sec:tradeoff:idle}
idle-energy correction's benefit is consistent across windows and
always positive (\GreenFaaS{} always beats v1).
\item Wait-Awhile and Lechowicz both consistently lose to FIFO across
all five windows ($-3.85\% \pm 1.99\%$ and $-5.56\% \pm 1.81\%$
respectively, both $\geq 2\sigma$ from zero on every window).
\end{itemize}

This is the strongest evidence we have that the paper's qualitative
findings are robust to real-world workload variability.
